\documentclass[11pt,twoside,a4paper]{article}
% http://www-h.eng.cam.ac.uk/help/tpl/textprocessing/latex_maths+pix/node6.html symboles de math
% http://fr.wikibooks.org/wiki/Programmation_LaTeX Programmation latex (wikibook)
%=========================== En-Tete =================================
%--- Insertion de paquetages (optionnel) ---
%--- Insertion de paquetages (optionnel) ---
\usepackage[french]{babel}   % pour dire que le texte est en fran{\'e}ais
\usepackage{a4}	             % pour la taille   
\usepackage[T1]{fontenc}     % pour les font postscript
\usepackage{epsfig}          % pour gerer les images
%\usepackage{psfig}
\usepackage{amsmath, amsthm} % tres bon mode mathematique
\usepackage{amsfonts,amssymb}% permet la definition des ensembles
\usepackage{float}           % pour le placement des figure
\usepackage{verbatim}


\usepackage{longtable} % pour les tableaux de plusieurs pages

\usepackage[table]{xcolor} % couleur de fond des cellules de tableaux

\usepackage{lastpage}

\usepackage{multirow}

\usepackage{multicol} % pour {\'e}crire dans certaines zones en colonnes : \begin{multicols}{nb colonnes}...\end{multicols}

%% https://texblog.org/2011/02/26/generating-dummy-textblindtext-with-latex-for-testing/
%% https://fr.sharelatex.com/learn/Multiple_columns
%% \usepackage{blindtext}
%% \usepackage{lipsum}
\usepackage{wrapfig}

% \usepackage[top=1.5cm, bottom=1.5cm, left=1.5cm, right=1.5cm]{geometry}
% gauche, haut, droite, bas, entete, ente2txt, pied, txt2pied
\usepackage{vmargin}
\setmarginsrb{1.0cm}{1.0cm}{1.0cm}{1.0cm}{15pt}{3pt}{57pt}{3pt}

\usepackage{lscape} % changement orientation page
\usepackage{pdflscape}
% --- style de page (pour les en-tete) ---
% \pagestyle{headings}
\pagestyle{empty}

% % % en-tete et pieds de page configurables : fancyhdr.sty

% http://www.trustonme.net/didactels/250.html

% http://ww3.ac-poitiers.fr/math/tex/pratique/entete/entete.htm
% http://www.ctan.org/tex-archive/macros/latex/contrib/fancyhdr/fancyhdr.pdf
% \usepackage{fancyhdr}
% \pagestyle{fancy}
% % \newcommand{\chaptermark}[1]{\markboth{#1}{}}
% % \newcommand{\sectionmark}[1]{\markright{\thesection\ #1}}
% \fancyhf{}
% \fancyhead[LE,RO]{\bfseries\thepage}
% \fancyhead[LO]{\bfseries\rightmark}
% \fancyhead[RE]{\bfseries\leftmark}
% \fancyfoot[LE]{\thepage /\pageref{LastPage} \hfill
	% TITLE
% \hfill \includegraphics[width=0.5cm]{img/logo_glider.png} }
% \fancyfoot[RO]{\includegraphics[width=0.5cm]{img/logo_glider.png} \hfill
	% TITLE
% \hfill \thepage /\pageref{LastPage}}
% \renewcommand{\headrulewidth}{0.5pt}
% \renewcommand{\footrulewidth}{0.5pt}
% \addtolength{\headheight}{0.5pt}
% \fancypagestyle{plain}{
	% \fancyhead{}
	% \renewcommand{\headrulewidth}{0pt}
% }

%--- Definitions de nouvelles commandes ---
\newcommand{\N}{\mathbb{N}} % les entiers naturels

%============================= Corps =================================
\begin{document}
\begin{landscape}

% \setcounter{page}{0}
% \thispagestyle{empty}

\begin{multicols}{2}
	\section*{Memento MAVEN}
	
	\begin{itemize}
		\item \textbf{POM} : Project Object Model (descripteur de projet maven)
		\item \textbf{Artefact} : fichier résultat du build d'un projet
		\item \textbf{Groupe} : Espace de nom permettant le regroupement d'artefacts
		\item \textbf{Version} : identification d'une version d'un artefact
		\item \textbf{Repository} : référentiel (dépôt) d'artefacts
		\item \textbf{Plugin} : composant enfichable dans le processus de build
		\item \textbf{Archétype} : modèle (prototype) de projet
		\item \textbf{Goal} : objectif du build, réalisé par un plugin
		%% \item[]
		\item \underline{\textbf{Phases : étapes clés de la construction du projet}}
		\begin{enumerate}
			\item \textsf{validate} : validation du projet, voir s'il est correct
			\item \textsf{compile} : compilation du projet
			\item \textsf{test} : compiler les tests unitaires et les passer
			\item \textsf{package} : mettre en forme le livrable. Souvent : faire un jar/un war
			\item \textsf{integration-test} : Tests qui doivent passer sur un environnement
			\item \textsf{verify} : Vérifier la qualité du livrable suivant une série de critères
			\item \textsf{install} : Déposer le package livrable sur le Repository Local
			\item \textsf{deploy} : Publier le package final sur un environnement partagé
		\end{enumerate}
%% validate
%% generate-sources
%% process-sources
%% generate-resources
%% process-resources
%% compile
%% process-test-sources
%% process-test-resources
%% test-compile
%% test
%% package
%% install
%% deploy
		\item hors cycle de vie : \textsf{clean}, \textsf{site}, \textsf{assembly:assembly}, \textsf{site-deploy}...
		%% \item[]
		\item \underline{\textbf{Convention plutôt que configuration}}
		\begin{itemize}
			\item \textsf{/src} : les sources du projet
			\item \textsf{/src/main} : code source et fichiers source principaux
			\item \textsf{/src/main/java} : code source
			\item \textsf{/src/main/resources} : fichiers de ressources (images, annexes...)
			\item \textsf{/src/main/webapp} : webapp du projet
			\item \textsf{/src/test} : fichiers de test
			\item \textsf{/src/test/java} : code source de test
			\item \textsf{/src/test/resources} : fichiers de ressources de test
			\item \textsf{/src/site} : informations sur le projet et/ou les rapports générés suite aux traitements effectués
			\item \textsf{/target} : fichiers résultat, les binaires (du code et des tests), les packages générés et les résultats des tests
		\end{itemize}
		\item[]
		\item \textbf{pom.xml} \texttt{http://maven.apache.org/maven-model/maven.html}
	\end{itemize}
	
	~\\
	
	\columnbreak
	
	\footnotesize
	\begin{verbatim}
<project xmlns="http://maven.apache.org/POM/4.0.0"
    xmlns:xsi="http://www.w3.org/2001/XMLSchema-instance"
    xsi:schemaLocation="http://maven.apache.org/POM/4.0.0
                        http://maven.apache.org/maven-v4_0_0.xsd">
    <modelVersion>4.0.0</modelVersion>
    
    <groupId>com.example.package</groupId>
    <artifactId>ApplicationName</artifactId>
    <packaging>jar</packaging>
    <version>1.0-SNAPSHOT</version>
    
    <name>Project Name</name>
    <url>http://www.example.org/</url>
    
    <dependencies>
        <!-- ... -->
        <dependency>
            <groupId>junit</groupId>
            <artifactId>junit</artifactId>
            <version>3.8.1</version>
            <scope>test</scope>
        </dependency>
        <!-- ... -->
    </dependencies>
    
    <properties>
        <maven.compiler.target>1.8</maven.compiler.target>
        <maven.compiler.source>1.8</maven.compiler.source>
        <build.name>BuildName</build.name>
        <!-- ... -->
    </properties>
    
    <build>
        <finalName>${build.name}</finalName>
        <plugins>
            <!-- ... -->
            <plugin>
                <groupId>org.apache.maven.plugins</groupId>
                <artifactId>maven-compiler-plugin</artifactId>
                <version>3.5.1</version>
                <configuration>
                    <source>${maven.compiler.source}</source>
                    <target>${maven.compiler.target}</target>
                </configuration>
            </plugin>
            <!-- ... -->
        </plugins>
    </build>
</project>
	\end{verbatim}
	\normalsize
	
	\vfill
	\columnbreak

	% \footnotesize
	% \begin{verbatim}
% <properties>
    % <maven.compiler.source>1.7</maven.compiler.source>
    % <maven.compiler.target>1.7</maven.compiler.target>

    % <spring-framework.version>4.0.6.RELEASE</spring-framework.version>
    % <slf4j.version>1.7.5</slf4j.version>

    % <project.build.sourceEncoding>UTF-8</project.build.sourceEncoding>
    % project.reporting.outputEncoding>UTF-8</project.reporting.outputEncoding>
% </properties>

% <dependencies>
    
    % <dependency>
        % <groupId>org.slf4j</groupId>
        % <artifactId>slf4j-api</artifactId>
        % <version>${slf4j.version}</version>
        % <scope>compile</scope>
    % </dependency>
    
% </dependencies>
	% \end{verbatim}
	% \normalsize
	
	
	...
\end{multicols}

\clearpage

\end{landscape}
\end{document}
